\section{Audit findings, verification, and integration status}
\label{rlc:sec:audit}

\subsection{Mathematical status and source overlap}
No new false equation was identified in the sampled canonical claims or
the incoming results used as context. Accordingly this delivery does
not manufacture an erratum. It contributes the proved reciprocal
identities above and records the exact normalization errors that would
arise from certain invalid simplifications.

The incoming bilinear article's same-direction integer-order master
\cite{BMH} and the present all-complex-order reciprocal master are
complementary. The former's arbitrary-order same-direction and
minimal-depth questions are not solved merely by replacing its kernel
with ours. The unequal-twist article \cite{UTJ} already establishes
substantial periodic mixed-jet results; our elementary rigidity and
partial-fraction arguments should be understood as their ordinary
real-line counterparts, not as a new general monodromy theorem.
The nonlinear-transport article \cite{NST} concerns coordinate-dependent
finite parts, whereas no finite-part transport is used in the present
convergent convolution family.

\begin{center}
\begin{tabular}{@{}p{0.54\textwidth}p{0.39\textwidth}@{}}
\toprule
Claim & Status in this delivery\\
\midrule
Reciprocal all-depth, all-complex-order closure & Proved, with weighted
holomorphy and finite formulas.\\[3pt]
Every mixed logarithmic moment & Proved by a terminating polynomial
recurrence; two-factor closed form supplied.\\[3pt]
Stieltjes products, harmonic lattice collapses, and anchored primitives
& Proved exact identities.\\[3pt]
Pure order-addition rigidity for unequal shifts & Proved within the
specified functional family.\\[3pt]
Canonical Gaussian $S_4$ & Already proved in the source; unchanged.\\[3pt]
Canonical Gaussian $S_6$ and revised $S_8$ & Not proved or disproved here.\\[3pt]
Historical priority for every displayed identity & Not claimed.\\[3pt]
Proof-assistant or interval certification & Not supplied.\\
\bottomrule
\end{tabular}
\end{center}

\subsection{Concrete safeguards and corrected intermediate work}
The following distinctions are necessary for a correct integration.

\paragraph{Retain entire resonance products.}
Use $\E_a(v)$ and the even-depth form
\eqref{rlc:eq:evencentral}, rather than evaluating a vanishing
Pochhammer symbol separately from a zeta pole. Formula
\eqref{rlc:eq:resonance} loses $1/3$ under that invalid operation.

\paragraph{Retain the first-moment factor.}
The correct first moment is $(s-t)\E_a(s+t+1)/2$.
An intermediate calculation during this work omitted the $1/2$.
The independent transform argument and the tilted-series coefficient
both detect the error. It has been corrected in the final formulas and
is retained as a deliberately failing negative control in the code.
This was an error in intermediate work, not an error attributed to the
repository.

\paragraph{Do not differentiate an integer-only theorem in its integer index.}
The unequal-shift partial fraction formula holds for specified positive
integer orders. Its finite number of terms does not itself extend
holomorphically when those integers are replaced by independent complex
orders. Use \eqref{rlc:eq:unequaltransform} or the appropriate Dirichlet
mixture for that operation. In contrast, the common-shift moment formula
\eqref{rlc:eq:momentclosed} is genuinely entire in its independent orders.

\paragraph{Keep branch-cancelled sums combined.}
At even depth and positive $\mu$, the individual Lerch terms in
\eqref{rlc:eq:lerchclosure} lie on a cut. Their jump cancels because
$p_r(0)=0$. A separated branch assignment that ignores that calculation
is not justified. At $0<\Rea a\le1/2$, use an admissible nonsymmetric
vertical line rather than the Gamma-weighted line $p=1/2+it$.

\paragraph{Keep remainders and derivative types distinct.}
The product in \eqref{rlc:eq:remainderidentity} is integrable even when
its separately expanded pieces are not. Also, $\partial_s\Li_s$,
$\partial_y\Li_s$, and $\Ima\Li_s$ are different operations.
The analytic Gaussian odd part \eqref{rlc:eq:Adef} is required at
complex order.

\subsection{Executed computational checks}
%% Generated from completed verification.json.
\newcommand{\ExactCount}{336}
\newcommand{\NumericCount}{91}
\newcommand{\NegativeCount}{3}
\newcommand{\WorkingDigits}{55}

The delivered script \src{code/verify.py} executed \ExactCount\ exact
symbolic assertions, \NumericCount\ floating-point diagnostics at
\WorkingDigits\ working decimal digits, and \NegativeCount\ negative
controls. The exact checks use rational polynomial algebra, with
$P=\pi^2$ treated as a formal variable. The numerical diagnostics use
independent contour and real-integral evaluations, not fitted integer
relations.

\begin{center}
\begin{tabular}{@{}lr@{}}
\toprule
Exact assertion family & Count\\
\midrule
Central products, parity, and cosecant differential equation &36\\
Two-factor pole-cancellation and reciprocity polynomials &32\\
Weighted-transform differential polynomials &9\\
Stirling versus Bell--harmonic derivatives &136\\
Unequal-shift partial fraction examples &4\\
Negative-order residue term &1\\
First-moment conservation &11\\
Independent all-depth versus two-factor moment reductions &25\\
Consecutive-shift complete-harmonic coefficients &80\\
Explicit triple and quadruple lattice values &2\\
\midrule
Total exact assertions &336\\
\bottomrule
\end{tabular}
\end{center}

The numerical families include complex and negative orders at depths
one through eight; shifts below $1/2$ on nonsymmetric contours;
direct convolutions of elementary logistic kernels; direct reciprocal
polylogarithm moments, including a noninteger-order example; nonzero
fugacity on both sides of the apparent positive cut; Stieltjes and
log-Gamma integrals; the two unequal-shift lattice examples; and the
arctangent and first harmonic-remainder integrals.

\textbf{Evidentiary limit.} These are floating-point diagnostics, not
rigorous enclosures. Working precision is not an asserted error bound.
The harmonic-remainder quadrature uses a finite logarithmic cutoff,
which is recorded in the script. The analytic proofs, not the small
residuals, establish the infinite and continuum assertions. The complete
machine-readable results, thresholds, package versions, and generated
coefficient tables are included in \src{data/}.

\subsection{Integration recommendation}
The package is additive. A natural home beneath the Mellin report is
\begin{center}\small
\src{continuations/reciprocal-lerch-convolutions/}
\end{center}
or a dedicated report path chosen by the maintainer.
All theorem and equation labels have the prefix \src{rlc:}.
The integration note contains a short proposed manuscript insertion,
a theorem index, the source pin, and explicit conjecture-status wording.
Neither the repository nor the user's persistent Library was changed.
